\section{The Four-Potential}
\label{sec:four_potential}

Relativity combines space and time into spacetime. The electromagnetic scalar and vector potentials must therefore be components of one Lorentz-covariant object. This object is the electromagnetic four-potential.

\subsection{Definition}

With

\[
x^{\mu}
=
(ct,\mathbf{x})
\]

and metric

\[
\eta_{\mu\nu}
=
\operatorname{diag}(1,-1,-1,-1),
\]

define the contravariant four-potential by

\[
\boxed{
A^{\mu}
=
\left(
\frac{\Phi}{c},
\mathbf{A}
\right)
}.
\]

In components,

\[
A^{0}
=
\frac{\Phi}{c},
\qquad
A^{i}=A_i^{\mathrm{(3D)}}.
\]

Lowering the index gives

\[
\boxed{
A_{\mu}
=
\left(
\frac{\Phi}{c},
-\mathbf{A}
\right)
}.
\]

Thus

\[
A_0=A^0,
\qquad
A_i=-A^i.
\]

The minus sign on the spatial components follows from the metric signature.

\subsection{Dimensional consistency}

The time component has units

\[
\left[
\frac{\Phi}{c}
\right]
=
\frac{\mathrm{V}}{\mathrm{m\,s^{-1}}}
=
\mathrm{V\,s\,m^{-1}}.
\]

These are the same units as the vector potential:

\[
[\mathbf{A}]
=
\mathrm{V\,s\,m^{-1}}.
\]

All four components can therefore belong to one geometric vector.

\subsection{Lorentz transformation law}

The four-potential transforms as a contravariant four-vector:

\[
\boxed{
A'^{\mu}
=
\Lambda^{\mu}{}_{\nu}A^{\nu}
}.
\]

This means that the scalar and vector potentials are not separately invariant. A Lorentz boost mixes them.

For a boost with velocity \(v\) in the positive \(x\)-direction,

\[
\beta
=
\frac{v}{c},
\qquad
\gamma
=
\frac{1}{\sqrt{1-\beta^2}}.
\]

The transformed components are

\[
A'^0
=
\gamma
\left(
A^0-
\beta A^1
\right),
\]

and

\[
A'^1
=
\gamma
\left(
A^1-
\beta A^0
\right).
\]

Therefore

\[
\boxed{
\Phi'
=
\gamma
\left(
\Phi-vA_x
\right)
}.
\]

The parallel vector-potential component transforms as

\[
\boxed{
A'_x
=
\gamma
\left(
A_x-
\frac{v}{c^2}\Phi
\right)
}.
\]

The transverse components are unchanged:

\[
A'_y=A_y,
\qquad
A'_z=A_z.
\]

A configuration that is described using only a scalar potential in one inertial frame can therefore require both scalar and vector potentials in another frame.

\subsection{Covariant gauge transformation}

The three-dimensional gauge transformation is

\[
\Phi'
=
\Phi-
\frac{\partial\chi}{\partial t},
\qquad
\mathbf{A}'
=
\mathbf{A}+
\nabla\chi.
\]

Using

\[
\partial^{\mu}
=
\left(
\frac{1}{c}
\frac{\partial}{\partial t},
-\nabla
\right),
\]

this becomes

\[
\boxed{
A'^{\mu}
=
A^{\mu}
-
\partial^{\mu}\chi
}.
\]

Equivalently, with lowered indices,

\[
\boxed{
A'_{\mu}
=
A_{\mu}
-
\partial_{\mu}\chi
}.
\]

The compact four-dimensional formula reproduces both three-dimensional transformations with the correct signs.

\subsection{The Lorenz gauge in covariant form}

The Lorenz gauge condition is

\[
\nabla\cdot\mathbf{A}
+
\frac{1}{c^2}
\frac{\partial\Phi}{\partial t}
=0.
\]

Compute the four-divergence:

\[
\begin{aligned}
\partial_{\mu}A^{\mu}
&=
\partial_0A^0
+
\partial_iA^i\\
&=
\frac{1}{c}
\frac{\partial}{\partial t}
\left(
\frac{\Phi}{c}
\right)
+
\nabla\cdot\mathbf{A}.
\end{aligned}
\]

Thus

\[
\boxed{
\partial_{\mu}A^{\mu}=0
}.
\]

The Lorenz gauge is manifestly Lorentz covariant because it is the vanishing of a four-scalar.

\subsection{Residual Lorenz-gauge freedom}

Under

\[
A'^{\mu}
=
A^{\mu}-\partial^{\mu}\chi,
\]

the divergence changes as

\[
\partial_{\mu}A'^{\mu}
=
\partial_{\mu}A^{\mu}
-
\partial_{\mu}\partial^{\mu}\chi.
\]

Therefore

\[
\partial_{\mu}A'^{\mu}
=
\partial_{\mu}A^{\mu}
-
\Box\chi.
\]

If the original potential satisfies the Lorenz gauge, the transformed potential also satisfies it when

\[
\boxed{
\Box\chi=0
}.
\]

\subsection{Four-current and source coupling}

Electric charge density and current density form the four-current

\[
\boxed{
J^{\mu}
=
\left(
c\rho,
\mathbf{J}
\right)
}.
\]

Lowering the index gives

\[
J_{\mu}
=
\left(
c\rho,
-\mathbf{J}
\right).
\]

The Lorentz scalar formed from the current and potential is

\[
\begin{aligned}
J_{\mu}A^{\mu}
&=
(c\rho)
\frac{\Phi}{c}
-
\mathbf{J}\cdot\mathbf{A}\\
&=
\rho\Phi
-
\mathbf{J}\cdot\mathbf{A}.
\end{aligned}
\]

The interaction density used in the electromagnetic action is

\[
\boxed{
\mathcal{L}_{\mathrm{int}}
=
-J_{\mu}A^{\mu}
=
-\rho\Phi
+
\mathbf{J}\cdot\mathbf{A}
}.
\]

This coupling will be varied in the next chapter.

\subsection{Gauge transformation of the source coupling}

Under a gauge transformation,

\[
\delta A^{\mu}
=
-\partial^{\mu}\chi.
\]

The interaction term changes by

\[
\delta\mathcal{L}_{\mathrm{int}}
=
J_{\mu}\partial^{\mu}\chi.
\]

Using the product rule,

\[
J_{\mu}\partial^{\mu}\chi
=
\partial^{\mu}
\left(
J_{\mu}\chi
\right)
-
\chi\partial^{\mu}J_{\mu}.
\]

If charge is locally conserved,

\[
\partial_{\mu}J^{\mu}=0,
\]

then the change is a total divergence. Gauge invariance of the source coupling is therefore tied to charge conservation.

\subsection{The potential norm is gauge dependent}

The contraction

\[
A_{\mu}A^{\mu}
=
\frac{\Phi^2}{c^2}
-
|\mathbf{A}|^2
\]

is a Lorentz scalar for a fixed four-potential, but it is not gauge invariant. Under

\[
A^{\mu}
\longrightarrow
A^{\mu}-\partial^{\mu}\chi,
\]

the value generally changes.

Lorentz invariance and gauge invariance are distinct requirements. A quantity can transform as a scalar under changes of inertial frame and still depend on the gauge representative.

\subsection{Plane-wave potential}

A vacuum plane-wave potential can be written as

\[
A^{\mu}(x)
=
\varepsilon^{\mu}
\cos(k_{\nu}x^{\nu}),
\]

where \(\varepsilon^{\mu}\) is the polarization four-vector.

The Lorenz gauge requires

\[
\partial_{\mu}A^{\mu}=0.
\]

Differentiating gives

\[
-k_{\mu}\varepsilon^{\mu}
\sin(k_{\nu}x^{\nu})=0.
\]

Therefore

\[
\boxed{
k_{\mu}\varepsilon^{\mu}=0
}.
\]

The polarization is four-dimensionally transverse to the wave vector. Residual gauge freedom can remove an additional unphysical component.

\subsection{Common mistakes}

Typical errors include:

\begin{enumerate}
    \item writing \(A_{\mu}=(\Phi/c,\mathbf{A})\) instead of changing the signs of the spatial components;
    \item treating \(\Phi\) and \(\mathbf{A}\) as separately Lorentz invariant;
    \item using \(A'^{\mu}=A^{\mu}+\partial^{\mu}\chi\) with the three-dimensional conventions adopted here;
    \item forgetting the factor \(1/c\) in the time component;
    \item confusing Lorentz covariance with gauge invariance;
    \item assuming that the Lorenz condition removes all gauge freedom.
\end{enumerate}

\subsection{What has been established}

The scalar and vector potentials combine into

\[
A^{\mu}
=
\left(
\frac{\Phi}{c},
\mathbf{A}
\right).
\]

They mix under Lorentz transformations and transform under gauge changes as

\[
A'^{\mu}
=
A^{\mu}-\partial^{\mu}\chi.
\]

The Lorenz gauge is the covariant condition \(\partial_{\mu}A^{\mu}=0\), and the source coupling is the Lorentz scalar \(-J_{\mu}A^{\mu}\). The next section constructs the gauge-invariant field tensor from derivatives of the four-potential.
